%% ============================================================
%% sections/a6-cost-estimate.tex
%% H. R. 9510 Bill v5.0 - APPENDIX A. COST ESTIMATE AND BUDGETARY DATA
%% (NON-OPERATIVE). The pre-introduction estimate in the CBO genre.
%% Rendered visuals: Table 7 (budget authority and outlays), Figure 8 (ASCII
%% waterfall), Table 8 (statutory user-fee structure), Table 9 (scorecard and
%% mandates summary).
%% ============================================================
% The appendices divider in main.tex supplies the page break for Appendix A.
\phantomsection\label{toc:a6}
\addcontentsline{toc}{section}{Appendix A. Cost Estimate and Budgetary Data (Non-Operative)}
\usccrosshead{APPENDIX A. COST ESTIMATE AND BUDGETARY DATA (NON-OPERATIVE)}

\uscprov{1}{\textit{This appendix is non-operative. It is a pre-introduction
estimate prepared in the genre of a Congressional Budget Office cost estimate;
it is not an official CBO product and is not endorsed by the CBO, the FDA,
HHS, or any Member of Congress. The formal estimate under section 402 of the
Congressional Budget Act of 1974 follows introduction and referral. Dollar
figures are illustrative unless tied to a cited statute or notice
\cite{cba-1974, cbo-paygo}.}}

\uscprov{1}{\textbf{Basis of the estimate.} H. R. 9510 sets one sequencing
rule (verification before generation), adds one financial-data record
(section 515D(k)), and directs one rulemaking, one records system, a
sequencing-inspection duty, and an annual publication (section 5). It creates
no entitlement, no new tax, and no new statutory fee; section 5(c) expressly
makes the filing of a verification record a non-fee-bearing event, and
section 5(d) authorizes - but does not provide - 58,000,000 dollars over
fiscal years 2027 through 2031. The estimate therefore has three parts: the
discretionary cost of the authorization (Table 7), the unchanged user-fee
structure the program rides on (Table 8), and the scorecard and mandate
determinations (Table 9) \cite{gao-budget-glossary, usc-379j}.}

\uscprov{1}{\textbf{Spending subject to appropriations.} Assuming
appropriation of the authorized amounts, estimated outlays follow the
spend-out profile typical of agency operating accounts: 60 percent in the
year of appropriation, 30 percent in the following year, and 10 percent in
the third year. Cumulative outlays within the window are 53.5 million
dollars of the 58.0 million dollars of budget authority; the remaining 4.5
million dollars outlays in fiscal years 2032 and 2033. Figure 8 draws the
ramp.}

{\footnotesize
\begin{xltabular}{\textwidth}{@{}Y R{1.35cm} R{1.35cm} R{1.35cm} R{1.35cm} R{1.35cm} R{1.5cm}@{}}
\hline\noalign{\vskip 0.04cm}
\textbf{(\$ millions, illustrative)} & \textbf{FY27} & \textbf{FY28} &
\textbf{FY29} & \textbf{FY30} & \textbf{FY31} & \textbf{FY27-31}\\
\hline\noalign{\vskip 0.04cm}
\endfirsthead
\hline\noalign{\vskip 0.04cm}
\textbf{(\$ millions, illustrative)} & \textbf{FY27} & \textbf{FY28} &
\textbf{FY29} & \textbf{FY30} & \textbf{FY31} & \textbf{FY27-31}\\
\hline\noalign{\vskip 0.04cm}
\endhead
\hline\noalign{\vskip 0.04cm}
\endlastfoot
Authorized budget authority (Table 6) & 18.0 & 12.0 & 10.0 & 9.0 & 9.0 &
58.0\\
\hline\noalign{\vskip 0.04cm}
Estimated outlays (60/30/10 spend-out) & 10.8 & 12.6 & 11.4 & 9.6 & 9.1 &
53.5\\
\hline\noalign{\vskip 0.04cm}
Outlays after the window (FY32-FY33) & - & - & - & - & - & 4.5\\
\hline
\end{xltabular}}
\vspace{-0.8cm}
\figcaption{Table 7. Estimated budget authority, fiscal years 2027 - 2031 (dollars in millions, illustrative). Each outlay cell applies the
60/30/10 spend-out to the Table 6 series; 53.5 plus the 4.5 after the window is the 58.0 authorized.}

\begin{asciifig}
ILLUSTRATIVE FIVE-YEAR PROFILE (authorized BA vs estimated outlays, $ millions)

         FY2027        FY2028        FY2029        FY2030        FY2031
  BA     18 |########  12 |######    10 |#####      9 |####       9 |####
            |              |             |             |             |
  OUT    11 |#####     13 |#######   11 |#####     10 |#####      9 |####
            |              |             |             |             |
  spend-out:  60% yr-1 + 30% yr-2 + 10% yr-3   (rulemaking + records build
  front-loaded); cumulative BA 58.0; outlays 53.5 in-window (rounded bars),
  remainder 4.5 in FY2032-FY2033
\end{asciifig}
\vspace{-0.75cm}
\figcaption{Figure 8. Budget authority to outlays (ASCII waterfall): the
authorized series beside the\\spend-out-driven outlay ramp, bars rounded to
whole millions; Table 7 carries the exact series.}

\uscprov{1}{\textbf{The user-fee structure the bill leaves unchanged.} The
medical-device user-fee program (sections 737 and 738; MDUFA V, fiscal years
2023 through 2027) prices each submission class as a statutory percentage of
the standard premarket approval fee, with dollar amounts set annually by
Federal Register notice. Table 8 states that structure; the bill's only
change to it is negative (no fee for a verification record filed alone), and
the sense of Congress in section 5(c)(2) points any verification-review
workload at the MDUFA VI negotiation \cite{usc-379i, usc-379j, pl-mdufa5,
fda-mdufa-rates}.}

{\footnotesize
\begin{xltabular}{\textwidth}{@{}Y R{5.4cm}@{}}
\hline\noalign{\vskip 0.04cm}
\textbf{Submission class (FD\&C \S~738; 21 U.S.C. \S~379j)} & \textbf{Statutory
percent of standard PMA fee}\\
\hline\noalign{\vskip 0.04cm}
\endfirsthead
\hline\noalign{\vskip 0.04cm}
\textbf{Submission class (FD\&C \S~738; 21 U.S.C. \S~379j)} & \textbf{Statutory
percent of standard PMA fee}\\
\hline\noalign{\vskip 0.04cm}
\endhead
\hline\noalign{\vskip 0.04cm}
\endlastfoot
Premarket approval application (PMA) & 100\\
\hline\noalign{\vskip 0.04cm}
Panel-track supplement & 75\\
\hline\noalign{\vskip 0.04cm}
De novo classification request & 30\\
\hline\noalign{\vskip 0.04cm}
180-day supplement & 15\\
\hline\noalign{\vskip 0.04cm}
Real-time supplement & 7\\
\hline\noalign{\vskip 0.04cm}
Annual fee for periodic reporting & 3.5\\
\hline\noalign{\vskip 0.04cm}
510(k) premarket notification & 3.4\\
\hline\noalign{\vskip 0.04cm}
30-day notice & 1.6\\
\hline\noalign{\vskip 0.04cm}
513(g) request for information & 1.35\\
\hline\noalign{\vskip 0.04cm}
Verification record filed alone (new \S~379j(a)(4)) & 0 (not fee-bearing)\\
\hline
\end{xltabular}}
\vspace{-0.78cm}
\figcaption{Table 8. The statutory device user-fee structure: each submission
class priced as a percent of\\the standard PMA fee under 21 U.S.C. \S~379j,
with the v5.0 addition - the verification record\\filed alone - priced at
zero. Annual dollar amounts are set by Federal Register notice.}

\uscprov{1}{\textbf{Sponsor compliance cost.} The private-sector cost of the
mandate is the verification run itself plus record-keeping: illustratively,
about 4,420 dollars per cleared interaction class (Appendix B, Table 10) and
about 45,000 dollars per investigation-year of documentation, attestation,
and cost-record administration, layered on validation work responsible
sponsors already perform. At those rates the mandate falls well below the
Unfunded Mandates Reform Act private-sector threshold (100,000,000 dollars in
1996 dollars, adjusted annually for inflation), and the bill imposes no
intergovernmental mandate; the State savings clause is preserved
\cite{usc-umra, usc-360k}.}

\uscprov{1}{\textbf{Scorecards and declarations.} Table 9 summarizes the
postures the operative text fixes: no direct spending or revenue effect, so
the Statutory PAYGO scorecard entry is expected to be zero under the section
5(e) clause; no net increase in mandatory spending, so the measure is CutGo
compatible; and no congressional earmark, limited tax benefit, or limited
tariff benefit as defined in clause 9 of Rule XXI - the authorization runs to
an agency for applicable implementation \cite{usc-paygo-2010,
house-rule-xxi}.}

{\footnotesize
\begin{xltabular}{\textwidth}{@{}L{4.6cm} Y@{}}
\hline\noalign{\vskip 0.04cm}
\textbf{Line item} & \textbf{Determination (pre-introduction, illustrative)}\\
\hline\noalign{\vskip 0.04cm}
\endfirsthead
\hline\noalign{\vskip 0.04cm}
\textbf{Line item} & \textbf{Determination (pre-introduction, illustrative)}\\
\hline\noalign{\vskip 0.04cm}
\endhead
\hline\noalign{\vskip 0.04cm}
\endlastfoot
Direct spending (mandatory) & None; no entitlement or mandatory program
created\\
\hline\noalign{\vskip 0.04cm}
Revenue & None; no tax, and the only fee change is a zero-fee carve-out\\
\hline\noalign{\vskip 0.04cm}
Statutory PAYGO scorecard & Expected zero; the section 5(e) clause governs\\
\hline\noalign{\vskip 0.04cm}
House CutGo (Rule XXI, clause 10) & Compatible; no net mandatory-spending
increase in either window\\
\hline\noalign{\vskip 0.04cm}
Discretionary cost & \$58.0M over five years
(illustrative); outlays per Table 7, subject to appropriations\\
\hline\noalign{\vskip 0.04cm}
Intergovern. mandate (UMRA) & None\\
\hline\noalign{\vskip 0.04cm}
Private-sector mandate (UMRA) & Exists; expected well below the threshold (per
the compliance-cost paragraph)\\
\hline\noalign{\vskip 0.04cm}
Earmarks, limited tax or tariff benefits (Rule XXI, cl. 9) & None (negative
declaration)\\
\hline\noalign{\vskip 0.04cm}
Formal CBO estimate & Follows introduction and referral to House Energy and
Commerce\\
\hline
\end{xltabular}}
\vspace{-0.8cm}
\figcaption{Table 9. Scorecard and mandates summary: the PAYGO, CutGo, UMRA,
and earmark determinations\\this pre-introduction estimate supports, each
traceable to the operative text it rests on.}
